\clearpage\pagestyle{otMain}\setlength{\otBlocEspacement}{2.5mm}\fontsize{11}{15}\selectfont\corrigesChapitre{15}{Angles et bissectrices}
\begin{blocLivre}{cor-15-1}\begin{corrige}{1}{Nommer un angle}\demarche{Le sommet s’écrit au milieu.}
L’angle se note $\widehat{AOB}$, ou $\widehat{BOA}$. Son sommet est $\res O$ ; ses côtés sont les demi-droites $\res{\text{[OA) et [OB)}}$.
\label{fig-cor-15-1}\figureGeo[corrige]{ang-1}
\end{corrige}\end{blocLivre}
\begin{blocLivre}{cor-15-2}\begin{corrige}{2}{Aigu, droit, obtus…}\demarche{Comparer chaque angle à un angle droit.}
a : $\res{\text{aigu}}$ ; b : $\res{\text{droit}}$ ; c : $\res{\text{obtus}}$ ; d : $\res{\text{plat}}$ ; e : $\res{\text{aigu}}$.
\label{fig-cor-15-2}\figureGeo[corrige]{ang-2}
\end{corrige}\end{blocLivre}
\begin{blocLivre}{cor-15-3}\begin{corrige}{3}{Lire le rapporteur}\demarche{Suivre la graduation qui part de 0.}
Le côté [OA) passe par le 0 de la graduation extérieure. On lit sur cette graduation : $\res{70^\circ}$. L’angle est aigu, ce qui confirme la lecture.
\label{fig-cor-15-3}\figureGeo[corrige]{ang-3}
\end{corrige}\end{blocLivre}
\begin{blocLivre}{cor-15-4}\begin{corrige}{4}{Avec l’autre graduation}\demarche{Contrôler : aigu ou obtus ?}
Le côté [OC) passe par le 0 de la graduation intérieure. On lit sur celle-ci : $\res{125^\circ}$, et non 55 : l’angle est obtus.
\label{fig-cor-15-4}\figureGeo[corrige]{ang-4}
\end{corrige}\end{blocLivre}
\begin{blocLivre}{cor-15-5}\begin{corrige}{5}{Estimer}\demarche{Comparer à l’angle droit et à l’angle plat.}
a est un peu plus grand qu’un angle droit : $\res{100^\circ}$. b est la moitié d’un angle droit : $\res{45^\circ}$. c est presque plat : $\res{160^\circ}$.
\label{fig-cor-15-5}\figureGeo[corrige]{ang-5}
\end{corrige}\end{blocLivre}
\begin{blocLivre}{cor-15-6}\begin{corrige}{6}{Construire un angle}\demarche{Placer le 0 sur le côté donné.}
Centre du rapporteur en A, 0 sur [Ax) ; on marque 65 et on trace [Ay), en vert. L’angle obtenu est aigu, comme prévu.
\label{fig-cor-15-6}\figureGeo[corrige]{ang-6}
\end{corrige}\end{blocLivre}
\begin{blocLivre}{cor-15-7}\begin{corrige}{7}{Adjacents ou non ?}\demarche{Vérifier les trois conditions.}
Figure 1 : $\res{\text{oui}}$, même sommet, côté commun, de part et d’autre. Figure 2 : $\res{\text{non}}$, l’un est à l’intérieur de l’autre. Figure 3 : $\res{\text{non}}$, les sommets sont différents.
\label{fig-cor-15-7}\figureGeo[corrige]{ang-7}
\end{corrige}\end{blocLivre}
\begin{blocLivre}{cor-15-8}\begin{corrige}{8}{Supplémentaires}\demarche{Compléter à 180 degrés.}
$\widehat{AOB}$ est plat : les deux angles sont supplémentaires. $180-128=\res{52^\circ}$.
\label{fig-cor-15-8}\figureGeo[corrige]{ang-8}
\end{corrige}\end{blocLivre}
\begin{blocLivre}{cor-15-9}\begin{corrige}{9}{Opposés par le sommet}\demarche{Combiner angles opposés et supplémentaires.}
L’angle opposé mesure aussi $\res{47^\circ}$. Les deux autres lui sont supplémentaires : $180-47=\res{133^\circ}$ chacun. Contrôle : $47+133+47+133=360$.
\label{fig-cor-15-9}\figureGeo[corrige]{ang-9}
\end{corrige}\end{blocLivre}
\begin{blocLivre}{cor-15-10}\begin{corrige}{10}{La bissectrice au rapporteur}\demarche{Diviser la mesure par deux.}
$84:2=\res{42^\circ}$. On marque 42 à partir de [Ox) et on trace la bissectrice, en vert.
\label{fig-cor-15-10}\figureGeo[corrige]{ang-10}
\end{corrige}\end{blocLivre}
\begin{blocLivre}{cor-15-11}\begin{corrige}{11}{Autour d’un point}\demarche{Utiliser l’angle plein.}
Un tour complet mesure $360^\circ$ : $360-110-95-80=\res{75^\circ}$.
\label{fig-cor-15-11}\figureGeo[corrige]{ang-11}
\end{corrige}\end{blocLivre}
\begin{blocLivre}{cor-15-12}\begin{corrige}{12}{Les aiguilles}\demarche{Compter des écarts de 30 degrés.}
Une heure correspond à $360:12=30^\circ$. À 3 h : $\res{90^\circ}$ ; à 4 h : $\res{120^\circ}$ ; à 5 h : $\res{150^\circ}$.
\label{fig-cor-15-12}\figureGeo[corrige]{ang-12}
\end{corrige}\end{blocLivre}
\begin{blocLivre}{cor-15-13}\begin{corrige}{13}{La bissectrice au compas}\demarche{Suivre les trois étapes de la construction.}
Arc de centre O, puis deux arcs de même rayon centrés sur les côtés : ils se coupent en P et [OP) est la bissectrice, en vert. Chaque moitié mesure $70:2=\res{35^\circ}$.
\label{fig-cor-15-13}\figureGeo[corrige]{ang-13}
\end{corrige}\end{blocLivre}
\begin{blocLivre}{cor-15-14}\begin{corrige}{14}{Deux bissectrices}\demarche{Additionner les deux moitiés.}
$\widehat{COB}=180-130=50^\circ$. Les moitiés mesurent $65^\circ$ et $25^\circ$. Les deux bissectrices forment un angle de $65+25=\res{90^\circ}$ : elles sont perpendiculaires.
\label{fig-cor-15-14}\figureGeo[corrige]{ang-14}
\end{corrige}\end{blocLivre}
\begin{blocLivre}{cor-15-15}\begin{corrige}{15}{Vrai ou faux}\demarche{Mesurer l’ouverture, pas les côtés.}
$\res{\text{Non}}$ : la longueur des côtés ne compte pas, seule l’ouverture compte. Au rapporteur, a mesure $40^\circ$ et b mesure $60^\circ$ : $\res{\text{b est plus grand}}$.
\label{fig-cor-15-15}\figureGeo[corrige]{ang-15}
\end{corrige}\end{blocLivre}
